\chapter{Ensembles y descomposición multiclase}
\section{Apuntes del Tema 2: Ensemble Models \& Multiclass Decomposition}

\subsection{1. Motivación y Métodos de Ensamble (\emph{Ensembles})}

\subsubsection{1.1 Filosofía del Ensamble}

\begin{itemize}
	\item \textbf{Limitaciones de un Modelo Individual:}
	\begin{itemize}
		\item Un modelo demasiado simple puede subajustar (\emph{underfit}), mostrando un alto sesgo (\emph{bias}) y bajo rendimiento.
		\item Un modelo excesivamente complejo puede sobreajustar (\emph{overfit}), memorizando el ruido de los datos de entrenamiento y mostrando alta varianza (\emph{variance}).
		\item Algunos modelos (como los árboles de decisión no podados) son inestables: pequeñas variaciones en los datos de entrenamiento generan cambios drásticos en la estructura del modelo.
	\end{itemize}
	
	\item \textbf{Idea Clave:} Combinar múltiples estimadores débiles (\emph{weak learners}) para construir un estimador fuerte (\emph{strong learner}) cuya predicción colectiva sea más precisa, estable y robusta que cualquiera de los modelos individuales por separado.
	
	\item \textbf{Diversidad:} Para que un ensamble funcione eficazmente, los estimadores individuales deben cometer errores de manera independiente o desacoplada (diversidad de hipótesis).
\end{itemize}

\subsection{2. Bagging y Random Forest}

\subsubsection{2.1 Bagging (\emph{Bootstrap Aggregating})}

\begin{itemize}
	\item \textbf{Objetivo:} Reducir la \textbf{varianza} de modelos inestables y con alta varianza (ej. árboles de decisión profundos) sin incrementar su sesgo.
	
	\item \textbf{Mecanismo Paso a Paso:}
	\begin{enumerate}
		\item \textbf{Muestreo Bootstrap:} Dado un dataset $\mathcal{D}$ de $N$ muestras, se generan $B$ subconjuntos $\mathcal{D}_b$ extrayendo $N$ muestras con reemplazo de $\mathcal{D}$.
		\item \textbf{Entrenamiento Independiente:} Se entrena un estimador $h_b(x)$ de forma completamente independiente sobre cada subconjunto $\mathcal{D}_b$.
		\item \textbf{Agregación:}
		\begin{itemize}
			\item \textbf{Clasificación:} Voto mayoritario:
			$$\hat{y} = \arg\max_{c} \sum_{b=1}^{B} \mathbb{I}(h_b(x) = c)$$
			\item \textbf{Regresión:} Promedio simple:
			$$\hat{y} = \frac{1}{B} \sum_{b=1}^{B} h_b(x)$$
		\end{itemize}
	\end{enumerate}
\end{itemize}

\subsubsection{2.2 Random Forest}

\begin{itemize}
	\item \textbf{Mejora sobre Bagging:} En el Bagging estándar, si existe una característica dominante, todos los árboles la elegirán en la raíz, resultando en un conjunto de árboles altamente correlacionados.
	
	\item \textbf{Decorrelación de Árboles:} Random Forest añade una segunda fuente de aleatoriedad: en cada división (\emph{split}) de cada árbol, se selecciona al azar un subconjunto reducido de $m$ características de las $d$ totales disponibles ($m \ll d$).
	\begin{itemize}
		\item Para clasificación: típicamente $m = \sqrt{d}$.
		\item Para regresión: típicamente $m = \frac{d}{3}$.
	\end{itemize}
	
	\item \textbf{Evaluación Out-Of-Bag (OOB):} Dado que cada muestra \emph{bootstrap} deja fuera aproximadamente un $36.8\%$ de los datos ($1 - 1/e \approx 0.632$ incluidos), estas muestras "no vistas" sirven como conjunto de validación implícito (\emph{OOB score}) sin requerir un split de validación adicional.
\end{itemize}

\subsection{3. Boosting: AdaBoost y Gradient Boosting (XGBoost)}

\subsubsection{3.1 AdaBoost (\emph{Adaptive Boosting})}

\begin{itemize}
	\item \textbf{Mecanismo:} Entrena estimadores de forma \textbf{secuencial}. Mantiene una distribución de pesos $D_t(i)$ sobre las muestras del dataset.
	
	\item \textbf{Algoritmo Paso a Paso:}
	\begin{enumerate}
		\item Inicializa los pesos de las muestras uniformemente: $D_1(i) = 1/N$.
		\item Para cada iteración $t = 1, \dots, T$:
		\begin{itemize}
			\item Entrena un estimador débil $h_t(x)$ minimizando el error ponderado:
			$$\epsilon_t = \sum_{i: y_i \neq h_t(x_i)} D_t(i)$$
			\item Calcula la importancia/confianza del estimador en la decisión final:
			$$\alpha_t = \frac{1}{2} \ln\left(\frac{1 - \epsilon_t}{\epsilon_t}\right)$$
			\item Actualiza los pesos de las muestras incrementando la importancia de las mal clasificadas:
			$$D_{t+1}(i) \propto D_t(i) \cdot \exp\left(-\alpha_t y_i h_t(x_i)\right)$$
			\item Renormaliza $D_{t+1}$ para que forme una distribución válida que sume 1.
		\end{itemize}
		\item \textbf{Predicción Final:} Voto ponderado:
		$$\hat{y} = \text{signo}\left(\sum_{t=1}^{T} \alpha_t h_t(x)\right)$$
	\end{enumerate}
\end{itemize}

\subsubsection{3.2 Gradient Boosting y XGBoost}

\begin{itemize}
	\item \textbf{Concepto Clave:} En lugar de reponderar muestras, cada nuevo árbol $h_t(x)$ se ajusta directamente para predecir los \textbf{residuos} (o el gradiente negativo de la función de pérdida) del ensamble acumulado $F_{t-1}(x)$.
	
	\item \textbf{Ajuste Aditivo:}
	$$F_t(x) = F_{t-1}(x) + \nu h_t(x)$$
	donde $\nu \in (0, 1]$ es el parámetro de tasa de aprendizaje (\emph{shrinkage}) que regula la contribución de cada nuevo árbol para evitar el sobreajuste.
	
	\item \textbf{Residuos para Error Cuadrático Medio (MSE):}
	$$r_{i, t-1} = -\left[ \frac{\partial \mathcal{L}(y_i, F(x_i))}{\partial F(x_i)} \right]_{F = F_{t-1}} = y_i - F_{t-1}(x_i)$$
	
	\item \textbf{XGBoost (\emph{Extreme Gradient Boosting}):} Extensión altamente eficiente que incluye:
	\begin{itemize}
		\item Expansión de Taylor de segundo orden (utiliza tanto el gradiente de primer orden $g_i$ como el Hessiano $h_i$).
		\item Regularización $L_1$ ($\alpha$) y $L_2$ ($\lambda$) sobre las hojas del árbol.
		\item Manejo automático e intrínseco de valores faltantes (\emph{sparse-aware split finding}).
	\end{itemize}
\end{itemize}

\subsection{4. Descomposición Multiclase: OVA y OVO}

Cuando un modelo de clasificación binario puro (ej. Regresión Logística, SVM) debe aplicarse a un problema de $K$ clases, se emplean estrategias de descomposición:

\begin{table}[htbp]
	\centering
	\small
	\begin{tabular}{|p{3cm}|p{2.5cm}|p{3.5cm}|p{4.5cm}|}
		\hline
		\textbf{Estrategia} & \textbf{Cantidad de Clasificadores} & \textbf{Configuración del Entrenamiento} & \textbf{Mecanismo de Predicción} \\ \hline
		\textbf{One-vs-All (OVA)} \emph{(One-vs-Rest)} & $K$ clasificadores & El clasificador $k$ se entrena con la clase $k$ como positiva (+1) y todas las demás como negativas (-1). & Se asigna la clase cuyo clasificador entregue la mayor probabilidad o puntuación (\emph{logit}): $\hat{y} = \arg\max_{k} f_k(x)$. \\ \hline
		\textbf{One-vs-One (OVO)} & $\frac{K(K-1)}{2}$ clasificadores & Se entrena un clasificador para cada par único de clases $(c_i, c_j)$ ignorando el resto. & Cada sub-clasificador emite un voto por su clase preferida. Gana la clase con \textbf{mayor número de votos acumulados}. \\ \hline
	\end{tabular}
\end{table}

\subsection{5. Práctica y Aplicación en Código (Scikit-Learn)}

\begin{lstlisting}[language=Python]
	import numpy as np
	from sklearn.datasets import make_classification
	from sklearn.model_selection import train_test_split
	from sklearn.ensemble import RandomForestClassifier, GradientBoostingClassifier, AdaBoostClassifier
	from sklearn.multiclass import OneVsRestClassifier, OneVsOneClassifier
	from sklearn.svm import SVC
	from sklearn.metrics import accuracy_score
	
	X, y = make_classification(
	n_samples=1000, 
	n_features=20, 
	n_informative=12, 
	n_classes=4, 
	random_state=42
	)
	
	X_train, X_test, y_train, y_test = train_test_split(
	X, y, test_size=0.2, random_state=42, stratify=y
	)
	
	def train_and_evaluate_ensembles(X_tr, y_tr, X_te, y_te):
	models = {
		"Random Forest": RandomForestClassifier(
		n_estimators=100, max_features="sqrt", random_state=42
		),
		"AdaBoost": AdaBoostClassifier(
		n_estimators=50, learning_rate=0.1, random_state=42
		),
		"Gradient Boosting": GradientBoostingClassifier(
		n_estimators=100, learning_rate=0.1, max_depth=3, random_state=42
		)
	}
	
	print("Evaluacion de Metodos de Ensamble")
	for name, model in models.items():
	model.fit(X_tr, y_tr)
	preds = model.predict(X_te)
	acc = accuracy_score(y_te, preds)
	print(f"Modelo: {name} | Test Accuracy: {acc:.4f}")
	
	def train_multiclass_decomposition(X_tr, y_tr, X_te, y_te):
	base_binary_clf = SVC(kernel='linear', C=1.0, random_state=42)
	
	ova_model = OneVsRestClassifier(base_binary_clf)
	ova_model.fit(X_tr, y_tr)
	ova_preds = ova_model.predict(X_te)
	
	ovo_model = OneVsOneClassifier(base_binary_clf)
	ovo_model.fit(X_tr, y_tr)
	ovo_preds = ovo_model.predict(X_te)
	
	print("Estrategias de Descomposicion Multiclase")
	print(f"Numero de estimadores OVA (K=4): {len(ova_model.estimators_)}")
	print(f"Accuracy OVA: {accuracy_score(y_te, ova_preds):.4f}")
	
	print(f"Numero de estimadores OVO: {len(ovo_model.estimators_)}")
	print(f"Accuracy OVO: {accuracy_score(y_te, ovo_preds):.4f}")
	
	if __name__ == "__main__":
	train_and_evaluate_ensembles(X_train, y_train, X_test, y_test)
	train_multiclass_decomposition(X_train, y_train, X_test, y_test)
\end{lstlisting}